\lecture{4.7}{Bootstrap Inference}
\label{lec:4-7}

\subsection{Why Variance Formulas Are Complex}

Point estimates alone are not enough. A causal analysis should also report uncertainty:
standard errors, confidence intervals, and sometimes hypothesis tests. For
g-computation and IPTW, uncertainty comes from several sources at once.

For g-computation, the estimator
\[
\psi_n^{\mathrm{gcomp}}(a)
=
\frac{1}{n}\sum_{i=1}^n \Qbar_n(a,W_i)
\]
depends on the empirical distribution of $W$ and on the fitted outcome regression
$\Qbar_n$. For IPTW, the estimator
\[
\psi_n^{\IPTW}(a)
=
\frac{1}{n}\sum_{i=1}^n
\frac{\ind(A_i=a)}{g_{n,a}(W_i)}Y_i
\]
depends on the empirical distribution of $(W,A,Y)$ and on the fitted propensity score
$g_{n,a}$. Treating $\Qbar_n$ or $g_n$ as fixed after fitting usually understates
uncertainty because it ignores the fact that the nuisance functions were estimated
from the same data.

\begin{keyidea}
The standard error must reflect the entire estimation procedure, not only the final
sample average. When nuisance functions are fit from data, their variability is part
of the variability of the causal effect estimator.
\end{keyidea}

\subsection{Nonparametric Bootstrap}

The nonparametric bootstrap approximates repeated sampling by resampling the observed
data. Let $\psi_n$ denote an estimator of a counterfactual mean or ATE. A bootstrap
replicate is constructed as follows:

\begin{enumerate}
\item Sample $n$ rows from the original dataset with replacement.
\item Refit all nuisance models in this bootstrap dataset.
\item Recompute the complete estimator in the bootstrap dataset.
\end{enumerate}

Repeating these steps $M$ times gives bootstrap estimates
\[
\psi_n^{\#,1},\psi_n^{\#,2},\ldots,\psi_n^{\#,M}.
\]
The symbol $\#$ is used here to distinguish bootstrap estimates from the original
estimate $\psi_n$.

\begin{figure}[H]
\centering
\begin{tikzpicture}[
  scale=.88,
  transform shape,
  node distance=.95cm,
  box/.style={draw, rounded corners=2pt, minimum width=3.0cm, minimum height=.78cm},
  arr/.style={-{Latex[length=2.5mm]}, thick}
]
\node[box] (data) {Original data};
\node[box, right=of data] (resample) {Resample rows\\ with replacement};
\node[box, right=of resample] (refit) {Refit nuisance\\ models};
\node[box, right=of refit] (estimate) {Recompute\\ estimator};
\draw[arr] (data) -- (resample);
\draw[arr] (resample) -- (refit);
\draw[arr] (refit) -- (estimate);
\draw[arr, rounded corners, dashed] (estimate.south) |- ++(0,-.8) -| node[below,pos=.25] {repeat $M$ times} (resample.south);
\end{tikzpicture}
\caption{Each bootstrap replicate repeats the full estimation procedure.}
\label{fig:ch4-bootstrap-flow}
\end{figure}

\subsection{Bootstrap Confidence Intervals}

A percentile bootstrap confidence interval uses quantiles of the bootstrap estimates.
For a nominal 95\% interval, let $\psi_n^\#(0.025)$ and $\psi_n^\#(0.975)$ be the
empirical 2.5th and 97.5th percentiles of
$\psi_n^{\#,1},\ldots,\psi_n^{\#,M}$. The interval is
\[
\left[
\psi_n^\#(0.025),\ \psi_n^\#(0.975)
\right].
\]
This interval is easy to compute and often adequate for introductory analyses, but it
can perform poorly in small samples, with highly skewed estimators, or when the
estimator is irregular.

\subsection{Bootstrap Standard Errors}

The bootstrap standard error is the sample standard deviation of the bootstrap
estimates:
\[
\hat\sigma^\#
=
\left[
\frac{1}{M-1}\sum_{m=1}^M
\left(\psi_n^{\#,m}-\bar\psi_n^\#\right)^2
\right]^{1/2},
\]
where
\[
\bar\psi_n^\#=\frac{1}{M}\sum_{m=1}^M \psi_n^{\#,m}.
\]
This standard error can be reported directly or used in Wald-style confidence
intervals of the form
\[
\psi_n\pm 1.96\,\hat\sigma^\#.
\]

\subsection{Wald-Style p-Values}

For testing
\[
H_0:\psi=0
\qquad\text{versus}\qquad
H_1:\psi\neq 0,
\]
a Wald-style test statistic is
\[
Z=\frac{\psi_n}{\hat\sigma^\#}.
\]
Using the standard normal approximation, the two-sided p-value is
\[
p=2\left\{1-\Phi(|Z|)\right\},
\]
where $\Phi$ is the cumulative distribution function of the standard normal
distribution.

This p-value inherits the assumptions of the large-sample approximation and the
bootstrap standard error. It should be interpreted as a statistical summary, not as a
measure of whether the causal assumptions are true.

\subsection{Bootstrap for G-Computation}

The key rule is that the outcome regression must be refit in every bootstrap sample.
If the fitted model from the original data is reused, the bootstrap will miss an
important component of uncertainty.

\begin{implementationbox}
Bootstrap for g-computation:
\begin{lstlisting}[language=R]
get_gcomp <- function(input_data) {
  fit_Q <- glm(Y ~ A + W1 + W2, data = input_data, family = binomial())

  data1 <- input_data
  data1$A <- 1
  data0 <- input_data
  data0$A <- 0

  Q1 <- predict(fit_Q, newdata = data1, type = "response")
  Q0 <- predict(fit_Q, newdata = data0, type = "response")

  psi1 <- mean(Q1)
  psi0 <- mean(Q0)
  c(psi0 = psi0, psi1 = psi1, ate = psi1 - psi0)
}

set.seed(1234)
M <- 1000
n <- nrow(dat)

boot_ate <- numeric(M)

for (m in seq_len(M)) {
  boot_ids <- sample(seq_len(n), size = n, replace = TRUE)
  boot_dat <- dat[boot_ids, ]
  boot_ate[m] <- get_gcomp(boot_dat)["ate"]
}

est <- get_gcomp(dat)["ate"]
se_boot <- sd(boot_ate)
ci_boot <- quantile(boot_ate, probs = c(0.025, 0.975))
p_value <- 2 * pnorm(-abs(est / se_boot))
\end{lstlisting}
\end{implementationbox}

The same bootstrap draws can be used to compute confidence intervals for
$\psi(0)$, $\psi(1)$, and the ATE if the function returns all three quantities.

\subsection{Bootstrap for IPTW}

For IPTW, the propensity score model must be refit in every bootstrap sample. The
weights in each bootstrap replicate should be computed from the bootstrap-fitted
propensity scores, not from the original fitted propensity scores.

\begin{implementationbox}
Bootstrap for IPTW:
\begin{lstlisting}[language=R]
get_iptw <- function(input_data) {
  fit_g <- glm(A ~ W1 + W2, data = input_data, family = binomial())

  g1 <- predict(fit_g, type = "response")
  g0 <- 1 - g1

  A <- input_data$A
  Y <- input_data$Y

  psi1 <- mean((A == 1) / g1 * Y)
  psi0 <- mean((A == 0) / g0 * Y)
  c(psi0 = psi0, psi1 = psi1, ate = psi1 - psi0)
}

set.seed(1234)
M <- 1000
n <- nrow(dat)

boot_ate <- numeric(M)

for (m in seq_len(M)) {
  boot_ids <- sample(seq_len(n), size = n, replace = TRUE)
  boot_dat <- dat[boot_ids, ]
  boot_ate[m] <- get_iptw(boot_dat)["ate"]
}

est <- get_iptw(dat)["ate"]
se_boot <- sd(boot_ate)
ci_boot <- quantile(boot_ate, probs = c(0.025, 0.975))
p_value <- 2 * pnorm(-abs(est / se_boot))
\end{lstlisting}
\end{implementationbox}

\begin{estimationbox}
Bootstrap checklist:
\begin{itemize}
\item Resample rows, not individual variables independently.
\item Refit every nuisance model inside each bootstrap sample.
\item Recompute the entire estimator from the resampled data.
\item Use enough bootstrap replicates for stable quantiles; $M=1000$ is a common
minimum for routine work.
\item Inspect failed model fits, separation warnings, and extreme weights, because
these are informative signs of instability.
\end{itemize}
\end{estimationbox}
